%% kapa: exam=1 round=26 year=2015 period=2 form=A subject=appraisal_law subject_ko=부동산관계법규 seq=34 number=34
%% source: Q-Net 감정평가사 기출문제 (cst00309) artlSeq=5138001 file=extracted/1008/5138001/제26회_감정평가사 1차_기출/제26회_감정평가사 1차_기출/2교시_감평_A형.pdf
%% conversion: kordoc 4.24.51 -> tools/gen1.py (auto); review=proofread-sonnet; flags=none
\begin{question}{34}
측량ㆍ수로조사 및 지적에 관한 법령상 지적공부와 등록사항의 연결이 바르지 않은 것은?
\begin{choices}
  \item 토지대장 - 지목과 면적
  \item 공유지연명부 - 소유권 지분
  \item 지적도 - 건축물 및 구조물 등의 위치
  \item 경계점좌표등록부 - 소유자와 부호도
  \item 대지권등록부 - 전유부분(專有部分)의 건물표시
\end{choices}
\end{question}
\answer{4}
